%% =====================================================================
\section{Limitations}
\label{sec:limitations}

\begin{itemize}
\item \emph{\Mseven wins on average, not on every puzzle, and its cost is high.} It is shorter than KMCoders' published solutions
  in mean length at every Santa size, but on some puzzles the published solution is still shorter (W/T/L in
  Table~\ref{tab:m7all}), and the Santa sizes have few puzzles (3--4 each from $N=6$ on), so the per-size margins are not
  significant differences. Each puzzle costs $k$ runs of KMCoders' program (minutes of CPU at $N=4$, over an hour at $N=8$--10 on
  four cores). We leave out the \Cu{19} and \Cu{33} Santa puzzles, where one copy takes hours. We have not yet compared it with an
  equal-compute portfolio of random seeds of the same program; the seed also changes the answer, and symmetries are one
  convenient, seed-free source of the same diversity.
\item \emph{\Mseven builds on KMCoders' code.} Its contribution is the portfolio, the shortened stage and the optimal central
  stage around their solver; the algorithmic core is theirs (\S\ref{sec:m7}). Their repository carries no licence file, so we
  release only a wrapper and a two-line patch, not their code.
\item \emph{Random-state samples.} The main tables use few states per size for large cubes (one each at $N=15$ and 20); the 20-state
  confidence intervals cover $N\le10$.
\item \emph{Still above Demaine's rate.} Every method's length divided by $N^2/\ln N$ still rises over the tested range, \FiveC's
  most slowly; KMCoders' lengths grow as $\approx4.4N^2$. None is shown to reach the $\Theta(N^2/\log N)$ rate in practice.
\item \emph{The counting bound is a lower bound.} It is valid for the radius around the goal, and we use it as an estimate of
  typical optimal distances; the Square calibration suggests the truth is 1.5--1.7$\times$ the bound in the worst case. It cannot
  certify that any single solution is near-optimal.
\item \emph{Exact search on the \Cu{3}} solves 12-move scrambles but no random state in 30 minutes on four cores. We did not
  train cube networks at CayleyPy's scale, so \Mfour and \FiveA are evaluated on the Square only.
\item \emph{Checker and sampling.} The checker enforces necessary conditions only, and random states come from long random walks
  rather than an exact uniform sampler.
\end{itemize}

%% =====================================================================
\section{Conclusion}

Decomposing Rubik's puzzles into clusters with $N$-invariant signatures lets one engine run eight solvers on the same instances,
from exact search to learned heuristics and a competition solver, for every size we can afford. On random states the winner at
large $N$ is structural: Demaine's idea of fixing many clusters with one shared move sequence, scheduled by net gain rather than
by unanimous need (\FiveC). On the Santa 2023 puzzles, the competition's own solver is far ahead of everything else, and the
symmetries of the cube are enough to turn it into a portfolio that is shorter on average than its published best-of-many
solutions at every Santa size up to $N=10$ (\Mseven). The
counting bound says where the remaining room is: little at $N=4$, and a factor that grows like $\log N$ at large $N$, because the
final stage of every practical method still fixes a few pieces per 3-cycle. Carrying \FiveC's bulk moves into the endgame of the
insertion method is the natural next step, and the place where the two lines of this work meet. Code, data and result files are
at \url{https://github.com/gauravrai1704/Rubik}.
